\setauthor{Markus Tran}
\section{Motivation und Problemstellung der Abo-Verwaltung}

Die zunehmende Verbreitung von Streaming-Diensten hat dazu geführt, dass viele Nutzer gleichzeitig mehrere Abonnements unterhalten. Dabei entsteht ein Problem, das in der Praxis häufig übersehen wird: Verschiedene Anbieter, wie etwa Netflix und Disney+, stellen teilweise identische Inhalte zur Verfügung. Dennoch werden beide Dienste parallel abonniert und abgerechnet, obwohl ein einziges Abonnement für den Zugriff auf den gewünschten Inhalt ausreichen würde. Dies führt zu unnötigen, wiederkehrenden Kosten für den Endnutzer.
Darüber hinaus fehlt es vielen Nutzern an einer zentralen Übersicht über ihre bestehenden Abonnements, was die bewusste Entscheidung für oder gegen einen bestimmten Anbieter erheblich erschwert. Die vorliegende Diplomarbeit greift diese Problemstellung auf und entwickelt eine Lösung, die Nutzern ermöglicht, fundierte Entscheidungen darüber zu treffen, welche Abonnements für sie tatsächlich einen Mehrwert bieten und welche gekündigt werden können.

\section{Zielsetzung der Diplomarbeit}
Ziel dieser Diplomarbeit ist die Konzeption und Entwicklung einer webbasierten Plattform zur zentralen Verwaltung von Streaming-Abonnements. Die Anwendung soll Nutzern ermöglichen, alle ihre aktiven Streaming-Dienste an einem einzigen Ort zu bündeln und jederzeit einen strukturierten Überblick über ihre Abonnements zu behalten.
Ein weiteres zentrales Ziel ist die gezielte Unterstützung bei der Inhaltssuche: Sucht ein Nutzer nach einem bestimmten Film oder einer Serie, soll die Plattform transparent anzeigen, bei welchen seiner bereits abonnierten Anbieter dieser Inhalt verfügbar ist. Damit wird nicht nur der Suchaufwand minimiert, sondern auch die Grundlage für eine bewusstere und kosteneffizientere Nutzung von Streaming-Diensten geschaffen.

\section{Nutzen der Anwendung für Endanwender}
Der unmittelbare Nutzen der entwickelten Anwendung für Endanwender liegt in der deutlichen Reduktion von Verwaltungsaufwand und unnötigen Ausgaben. Durch eine übersichtliche und intuitiv bedienbare Benutzeroberfläche erhalten Nutzer auf einen Blick eine vollständige Darstellung ihrer aktiven Abonnements sowie deren Kosten.
Zusätzlich profitieren Nutzer von einer intelligenten Inhaltssuche, die vorhandene Abonnements berücksichtigt und den günstigsten bzw. bereits bezahlten Anbieter für einen gesuchten Inhalt hervorhebt. Insgesamt trägt die Anwendung dazu bei, das Management mehrerer Streaming-Dienste erheblich zu vereinfachen und die finanzielle Transparenz für den Nutzer zu erhöhen.

\section{Aufbau der Diplomarbeit}
Die vorliegende Arbeit gliedert sich in mehrere thematische Abschnitte. Nach der Einleitung und der Analyse der Ausgangssituation werden die technischen Grundlagen und die eingesetzten Technologien vorgestellt.
Die Implementierung basiert auf einem modernen Technologie-Stack: Als Backend-Framework wird Quarkus eingesetzt, während das Frontend mit Angular realisiert wird. Für die Authentifizierung und das Identitätsmanagement kommt Keycloak als externe Bibliothek zum Einsatz. Die Bereitstellung von Film- und Seriendaten erfolgt über die The Movie Database (TMDb) API, während eine relationale Datenbank die persistente Speicherung und Verwaltung der Nutzerdaten übernimmt.
Abschließend werden die Ergebnisse evaluiert sowie ein Ausblick auf mögliche Weiterentwicklungen gegeben.